\documentclass[12pt]{article}
\usepackage{amsmath,amssymb,amsthm}
\usepackage{mathtools}
\usepackage{geometry}
\geometry{margin=1in}

% Theorem environments
\newtheorem{definition}{Definition}[section]
\newtheorem{lemma}[definition]{Lemma}
\newtheorem{theorem}[definition]{Theorem}
\newtheorem{corollary}[definition]{Corollary}
\newtheorem{proposition}[definition]{Proposition}
\newtheorem{remark}[definition]{Remark}

% Custom commands
\newcommand{\B}{\mathbb{B}}
\newcommand{\C}{\mathbb{C}}
\newcommand{\R}{\mathbb{R}}
\renewcommand{\H}{\mathbb{H}}
\newcommand{\Tr}{\mathrm{Tr}}
\newcommand{\re}{\mathrm{Re}}
\newcommand{\im}{\mathrm{Im}}

\title{Step 1: Metric Variation of the UBT Matter Action\\
\large Derivation of the Stress-Energy Tensor $T_{\mu\nu}$ from $S[\Theta]$}
\author{UBT Theory Development}
\date{2026}

\begin{document}

\maketitle

\begin{abstract}
We derive the stress-energy tensor $T_{\mu\nu}$ of the UBT matter sector via the
Hilbert (metric variation) prescription.  Starting from the complete matter action
$S_{\mathrm{matter}}[\Theta, g] = S_{\mathrm{kin}} + S_{\mathrm{pot}} +
S_{\mathrm{gauge}}$ defined in \texttt{appendix\_AA\_theta\_action.tex}, we
compute $\delta S_{\mathrm{matter}}/\delta g^{\mu\nu}$ for each term explicitly,
step by step.  The results close the gap in
\texttt{appendix\_R\_GR\_equivalence.tex} where $T_{\mu\nu} = \mathrm{Re}(\mathcal{T}_{\mu\nu})$
is stated without derivation.
\end{abstract}

\tableofcontents

% ------------------------------------------------------------
\section{Setup and Conventions}

\subsection{Metric Signature and Notation}

We work in Lorentzian signature $(-,+,+,+)$ throughout:
$g_{\mu\nu} = \mathrm{diag}(-1,+1,+1,+1)$ in flat space.
Indices are raised and lowered with $g^{\mu\nu}$ and $g_{\mu\nu}$.
We write $g := \det(g_{\mu\nu}) < 0$ and $\sqrt{-g} > 0$.

The covariant derivative acting on $\Theta$ is:
\begin{equation}
  \nabla_\mu = \partial_\mu + \Omega_\mu + i g_{\mathrm{YM}} A_\mu,
\end{equation}
where $\Omega_\mu$ is the spin connection and $A_\mu$ is the gauge connection.

\subsection{Hilbert Stress-Energy Tensor}

\begin{definition}[Hilbert stress-energy tensor]
\label{def:hilbert}
Given a matter action $S_{\mathrm{matter}}[\Theta, g]$, the
\emph{Hilbert (metric variation) stress-energy tensor} is:
\begin{equation}
  \label{eq:hilbert}
  T_{\mu\nu} := -\frac{2}{\sqrt{-g}}\,
    \frac{\delta S_{\mathrm{matter}}}{\delta g^{\mu\nu}}.
\end{equation}
\end{definition}

\begin{remark}
  The sign in \eqref{eq:hilbert} is chosen so that the matter sector in
  Einstein's equations contributes $+8\pi G\,T_{\mu\nu}$ on the right-hand side.
  The convention follows Wald~\cite{wald} and Carroll~\cite{carroll}.
\end{remark}

\subsection{Matter Action}

The UBT matter action (from \texttt{appendix\_AA\_theta\_action.tex},
Section~3) read, in the effective 4D Lorentzian projection:
\begin{align}
  S_{\mathrm{kin}}   &= -\tfrac{1}{2}\int d^4x\,\sqrt{-g}\;
    g^{\alpha\beta}\,\Tr\!\bigl[(\nabla_\alpha\Theta)^\dagger\,
    \nabla_\beta\Theta\bigr], \label{eq:skin} \\
  S_{\mathrm{pot}}   &= -\int d^4x\,\sqrt{-g}\; V(\Theta),
    \label{eq:spot} \\
  S_{\mathrm{gauge}} &= -\frac{K}{4}\int d^4x\,\sqrt{-g}\;
    \Tr[F_{\mu\nu}F^{\mu\nu}].
    \label{eq:sgauge}
\end{align}
Here $K = K_{\mathrm{gauge}} > 0$ is the gauge kinetic normalisation
($K = 1$ for the Abelian case), and $F_{\mu\nu} = \partial_\mu A_\nu - \partial_\nu A_\mu
+ i g_{\mathrm{YM}}[A_\mu, A_\nu]$.

\begin{remark}[Sign of $S_{\mathrm{kin}}$]
  The minus sign in \eqref{eq:skin} is required by the $(-,+,+,+)$ convention.
  In this signature $g^{\alpha\beta}(\nabla_\alpha\Theta)^\dagger\nabla_\beta\Theta
  = -({\partial_t\Theta})^2 + (\nabla\Theta)^2$,
  so a minus sign in the action yields positive kinetic energy density
  $\tfrac{1}{2}(\partial_t\Theta)^2 + \tfrac{1}{2}(\nabla\Theta)^2 > 0$.
\end{remark}

\subsection{Standard Variational Identities}

\begin{lemma}[Metric variation identities]
\label{lem:var_ids}
Under a variation $g^{\mu\nu} \to g^{\mu\nu} + \delta g^{\mu\nu}$:
\begin{align}
  \delta\sqrt{-g}
    &= -\tfrac{1}{2}\sqrt{-g}\; g_{\mu\nu}\,\delta g^{\mu\nu},
    \label{eq:var_sqrtg} \\
  \delta g_{\alpha\beta}
    &= -g_{\alpha\mu}\,g_{\beta\nu}\,\delta g^{\mu\nu},
    \label{eq:var_lower} \\
  \delta\!\bigl(\sqrt{-g}\,g^{\alpha\beta}\bigr)
    &= \sqrt{-g}\Bigl[\delta^{(\alpha}_\mu\delta^{\beta)}_\nu
       -\tfrac{1}{2}g_{\mu\nu}g^{\alpha\beta}\Bigr]\delta g^{\mu\nu}.
    \label{eq:var_combined}
\end{align}
\end{lemma}

\begin{proof}
  Identity~\eqref{eq:var_sqrtg}: from $\delta\ln|\det g| = g^{\mu\nu}\delta g_{\mu\nu}
  = -g_{\mu\nu}\delta g^{\mu\nu}$ (using $g^{\mu\alpha}g_{\alpha\nu}=\delta^\mu_\nu$).
  Identity~\eqref{eq:var_lower}: differentiate $g^{\mu\alpha}g_{\alpha\nu}=\delta^\mu_\nu$.
  Identity~\eqref{eq:var_combined}: combine the first two with the product rule
  and symmetrise in $(\mu,\nu)$.
\end{proof}

% ------------------------------------------------------------
\section{Variation of the Kinetic Term}
\label{sec:kin}

\subsection{Explicit Computation}

We write $S_{\mathrm{kin}} = -\tfrac{1}{2}\int d^4x\,
\sqrt{-g}\,g^{\alpha\beta}\,K_{\alpha\beta}$, where
$K_{\alpha\beta} := \Tr[(\nabla_\alpha\Theta)^\dagger\nabla_\beta\Theta]$.

Varying with respect to $g^{\mu\nu}$:
\begin{equation}
  \delta S_{\mathrm{kin}}
  = -\tfrac{1}{2}\int d^4x\,\Bigl[
      \delta(\sqrt{-g}\,g^{\alpha\beta})\,K_{\alpha\beta}
    + \sqrt{-g}\,g^{\alpha\beta}\,\delta K_{\alpha\beta}
    \Bigr].
\end{equation}

\paragraph{First piece (explicit metric dependence).}
Using identity~\eqref{eq:var_combined}:
\begin{equation}
  \delta(\sqrt{-g}\,g^{\alpha\beta})\,K_{\alpha\beta}
  = \sqrt{-g}\Bigl[K_{\mu\nu} - \tfrac{1}{2}g_{\mu\nu}\,g^{\alpha\beta}K_{\alpha\beta}
    \Bigr]\delta g^{\mu\nu}.
\end{equation}

\paragraph{Second piece (connection variation).}
The connection $\Omega_\mu$ depends on $g_{\mu\nu}$ through the tetrad.
Under metric variation, $\delta\Omega_\mu \sim \delta\Gamma^\lambda_{\mu\nu}$,
producing terms of the form
$\sqrt{-g}\,g^{\alpha\beta}\Tr[(\delta\nabla_\alpha\Theta)^\dagger\nabla_\beta\Theta]$.
Integrating by parts converts these into bulk terms proportional to
$(\nabla^\mu\nabla_\mu\Theta)\,\delta\Theta$ plus boundary terms.
When $\Theta$ satisfies the Euler--Lagrange equation
$\nabla^\mu\nabla_\mu\Theta = \partial V/\partial\Theta^\dagger$
(Theorem~4.1 of \texttt{appendix\_AA}), the connection variation contributes
only boundary terms that vanish for variations with compact support.
We therefore \emph{discard} this piece \textbf{on-shell}.

\begin{proposition}[Kinetic stress-energy tensor]
\label{prop:T_kin}
For $\Theta$ satisfying the UBT field equations:
\begin{equation}
  \frac{\delta S_{\mathrm{kin}}}{\delta g^{\mu\nu}}
  = \tfrac{1}{2}\sqrt{-g}\Bigl[
      K_{\mu\nu} - \tfrac{1}{2}g_{\mu\nu}\,K^\alpha{}_\alpha
    \Bigr],
\end{equation}
where $K^\alpha{}_\alpha = g^{\alpha\beta}K_{\alpha\beta}
= \Tr[(\nabla_\alpha\Theta)^\dagger\nabla^\alpha\Theta]$.
Hence:
\begin{equation}
  \label{eq:T_kin}
  \boxed{
    T^{(\mathrm{kin})}_{\mu\nu}
    = \Tr\!\bigl[(\nabla_\mu\Theta)^\dagger\nabla_\nu\Theta\bigr]
      -\tfrac{1}{2}g_{\mu\nu}\,
       \Tr\!\bigl[(\nabla_\alpha\Theta)^\dagger\nabla^\alpha\Theta\bigr].
  }
\end{equation}
\end{proposition}

\begin{proof}
  Direct application of Definition~\ref{def:hilbert} to the first piece:
  $T^{(\mathrm{kin})}_{\mu\nu} = -(2/\sqrt{-g})\times
  (-\tfrac{1}{2})\times\tfrac{1}{2}\sqrt{-g}[K_{\mu\nu} -
  \tfrac{1}{2}g_{\mu\nu}K^\alpha{}_\alpha]
  = K_{\mu\nu} - \tfrac{1}{2}g_{\mu\nu}K^\alpha{}_\alpha$.
\end{proof}

\begin{remark}[Scalar field check]
  For a real scalar $\phi$ (setting $\Tr[\cdot] \to \tfrac{1}{2}\cdot$):
  $T^{(\mathrm{kin})}_{\mu\nu} = \partial_\mu\phi\,\partial_\nu\phi
  - \tfrac{1}{2}g_{\mu\nu}(\partial\phi)^2$,
  the standard Klein--Gordon kinetic stress-energy tensor.~$\checkmark$
\end{remark}

% ------------------------------------------------------------
\section{Variation of the Potential Term}
\label{sec:pot}

\begin{proposition}[Potential stress-energy tensor]
\label{prop:T_pot}
The potential $V(\Theta)$ depends on $g^{\mu\nu}$ only through the
integration measure $\sqrt{-g}$ (it is a scalar built from $\Theta$ with
no explicit metric contractions).  Therefore:
\begin{equation}
  \frac{\delta S_{\mathrm{pot}}}{\delta g^{\mu\nu}}
  = -\int d^4x\,\frac{\delta\sqrt{-g}}{\delta g^{\mu\nu}}\,V(\Theta)
  = \tfrac{1}{2}\sqrt{-g}\;g_{\mu\nu}\,V(\Theta).
\end{equation}
The potential contribution to the stress-energy tensor is:
\begin{equation}
  \label{eq:T_pot}
  \boxed{T^{(\mathrm{pot})}_{\mu\nu} = -g_{\mu\nu}\,V(\Theta).}
\end{equation}
\end{proposition}

\begin{remark}[Vacuum energy interpretation]
  $T^{(\mathrm{pot})}_{\mu\nu} = -g_{\mu\nu}V$ has the form of a
  cosmological constant.  For the Mexican-hat potential
  $V = \tfrac{\lambda}{4}(|\Theta|^2-v^2)^2$, at the vacuum
  $|\Theta|=v$ we have $V=0$ and $T^{(\mathrm{pot})}_{\mu\nu}=0$.
\end{remark}

% ------------------------------------------------------------
\section{Variation of the Gauge Kinetic Term}
\label{sec:gauge}

\subsection{Decomposing the Metric Dependence}

With $F^{\mu\nu} = g^{\mu\alpha}g^{\nu\beta}F_{\alpha\beta}$, the
integrand is $\sqrt{-g}\,g^{\mu\alpha}g^{\nu\beta}F_{\alpha\beta}F_{\mu\nu}$.
The metric appears in three independent factors: $\sqrt{-g}$, the first
$g^{\mu\alpha}$, and the second $g^{\nu\beta}$.

\subsection{Explicit Computation}

\begin{proposition}[Gauge kinetic stress-energy tensor]
\label{prop:T_gauge}
Varying $S_{\mathrm{gauge}}$ term by term with respect to $g^{\rho\sigma}$:

\paragraph{Variation of $\sqrt{-g}$:}
\begin{equation}
  -\frac{K}{4}\Bigl(-\tfrac{1}{2}\sqrt{-g}\,g_{\rho\sigma}\Bigr)
  F_{\mu\nu}F^{\mu\nu}
  = \frac{K}{8}\sqrt{-g}\;g_{\rho\sigma}\Tr[F_{\mu\nu}F^{\mu\nu}].
\end{equation}

\paragraph{Variation of the first $g^{\mu\alpha}$:}
\begin{equation}
  -\frac{K}{4}\sqrt{-g}\;\delta^{(\mu}_\rho\delta^{\alpha)}_\sigma
  \,g^{\nu\beta}\Tr[F_{\alpha\beta}F_{\mu\nu}]
  = -\frac{K}{4}\sqrt{-g}\;\Tr[F_{\rho\beta}F_\sigma{}^\beta
    + F_{\sigma\beta}F_\rho{}^\beta].
\end{equation}

\paragraph{Variation of the second $g^{\nu\beta}$:}
By the antisymmetry $F_{\mu\nu} = -F_{\nu\mu}$,
this is identical to the previous piece.

\paragraph{Total:}
\begin{equation}
  \frac{\delta S_{\mathrm{gauge}}}{\delta g^{\rho\sigma}}
  = \sqrt{-g}\,\frac{K}{4}\Bigl[
      \tfrac{1}{2}g_{\rho\sigma}\Tr[F_{\mu\nu}F^{\mu\nu}]
      - 2\Tr[F_{\rho\alpha}F_\sigma{}^{\alpha}]
    \Bigr].
\end{equation}
Applying Definition~\ref{def:hilbert}:
\begin{equation}
  \label{eq:T_gauge}
  \boxed{
    T^{(\mathrm{gauge})}_{\mu\nu}
    = K\Bigl[
        \Tr\!\bigl[F_{\mu\alpha}F_\nu{}^{\alpha}\bigr]
        -\tfrac{1}{4}g_{\mu\nu}\Tr\!\bigl[F_{\alpha\beta}F^{\alpha\beta}\bigr]
      \Bigr].
  }
\end{equation}
\end{proposition}

\begin{remark}[Electromagnetic limit]
  For $K=1$ and Abelian $F_{\mu\nu}$ (with $\Tr[\cdot]\to\cdot$):
  $T^{(\mathrm{EM})}_{\mu\nu} = F_{\mu\alpha}F_\nu{}^\alpha
  - \tfrac{1}{4}g_{\mu\nu}F_{\alpha\beta}F^{\alpha\beta}$,
  which is the standard Maxwell stress-energy tensor.
  This is traceless in $d=4$: $g^{\mu\nu}T^{(\mathrm{EM})}_{\mu\nu}=0$.~$\checkmark$
\end{remark}

% ------------------------------------------------------------
\section{Summary of Results}
\label{sec:summary}

\begin{theorem}[Metric variation of UBT matter action]
\label{thm:T_matter}
The complete UBT matter stress-energy tensor, obtained by the Hilbert
prescription from $S_{\mathrm{matter}} = S_{\mathrm{kin}}
+ S_{\mathrm{pot}} + S_{\mathrm{gauge}}$, is:
\begin{equation}
  \label{eq:T_total}
  \boxed{
    T_{\mu\nu}
    \;=\;
    \underbrace{\Tr[(\nabla_\mu\Theta)^\dagger\nabla_\nu\Theta]
      -\tfrac{1}{2}g_{\mu\nu}\Tr[(\nabla_\alpha\Theta)^\dagger\nabla^\alpha\Theta]}_{%
      T^{(\mathrm{kin})}_{\mu\nu}}
    \;-\;
    \underbrace{g_{\mu\nu}V(\Theta)}_{%
      -T^{(\mathrm{pot})}_{\mu\nu}}
    \;+\;
    K\underbrace{\Bigl[\Tr[F_{\mu\alpha}F_\nu{}^\alpha]
      -\tfrac{1}{4}g_{\mu\nu}\Tr[F_{\alpha\beta}F^{\alpha\beta}]\Bigr]}_{%
      T^{(\mathrm{gauge})}_{\mu\nu}/K}.
  }
\end{equation}
This derivation is on-shell (valid when $\Theta$ satisfies the Euler--Lagrange
equations of Appendix~AA, Theorem~4.1).
\end{theorem}

\begin{remark}[Status]
  Equation~\eqref{eq:T_total} is derived purely from the existing action in
  \texttt{appendix\_AA\_theta\_action.tex} with no new free parameters.
  No new physics is introduced.
  The biquaternionic generalisation $\mathcal{T}_{\mu\nu}$ is constructed in
  \texttt{step2\_biquat\_stress\_energy.tex}.
  The Einstein equations $G_{\mu\nu} = 8\pi G\,T_{\mu\nu}$ are derived in
  \texttt{step3\_einstein\_with\_matter.tex}.
\end{remark}

\begin{thebibliography}{99}
\bibitem{wald}    Wald, R.~M. (1984). \textit{General Relativity}.
                  University of Chicago Press.
\bibitem{carroll} Carroll, S.~M. (2004). \textit{Spacetime and Geometry}.
                  Addison--Wesley.
\bibitem{ps}      Peskin, M.~E.\ \& Schroeder, D.~V. (1995).
                  \textit{An Introduction to Quantum Field Theory}. Westview Press.
\bibitem{appendix_AA}
                  UBT Theory Development.
                  \texttt{appendix\_AA\_theta\_action.tex}, 2025.
\end{thebibliography}

\end{document}
